\scanpage{8}\mspage{\handwritten{4}}
\hfill Bures le 12.8.19\deleted{\unsure{4|8|6}}64\\
Mon cher Serre,\\

Je commence à avoir de faibles lumières sur la façon de faire \handwritten{d}es VA avec les cycles al\handwritten{g}ébriquement équivalents à zéro d'une variété projective non singulière $X$ (sur un corps alg clos $k$). Si $X$ est connexe de dimension $n$, je sais associer à chaque entier $i$' \deleted{\unsure{ç|c}}compris entre $\deleted{\unsure{0|1}}1$ et $n$ une VA $J^i(X)$, jouant le rôle d'une "jacobienne intermédiaire" au sens de Weil, correspondant à des classes de cycles de codimension $i$ (pour une relation d'équivalence que je ne sais pas bien identifier pour l'instant, mais qui mérite certainement d'être appelée "équivalence d'Albanese"). Je suis moralement sûr que ce sont"les bonnes", bien que je n'ai pas prouvé grand-chose dessus pour l'instant; par contre, j'ai des conjectures en pagaille. Je te signale seulement que $J^1=\deleted{\unsure{B|P}}\operatorname{Pic}^{\unsure{o|0}}$ et $J^n=\operatorname{Alb}^{o}$, que $J^i$ et $J^{n+1-i}$ sont canoniquement duales l'une de l'autre, que $\dim J^i\handwritten{\leq\frac{1}{2}}b_{2i-1}$ (nombre de Bett\deleted{\unsure{a|e}}i) - de façon plus précise $T_{\unsure{\ell|l}}(J^i)$ est un quotient d'un sous-module de $H^{2i-1}(X,\underline{Z}_{\ell}(i))$ (et quand on sera plus savant, on saura prouver que c'est même un sous-module dudit), en caractéristique $0$ on peut également remplacer la cohomologie de Weil par la cohomologie de Hodge $H^i(X,\handwritten{\underline{\Omega}}_X^{i-1})$ ...

% slanted; two separate notes next to the Hodge formula
\margin{\handwritten{conjecturalement $\unsure{\underline{t}_{J^i}|t_{J^i}|\mathcal{Z}^{i}}\unsure{\subseteq|\subset} H^i(X,\underline{\Omega}_X^{i-1})$}}
\margin{\handwritten{\deleted{\unsure{$H^1(J^i)$|$H^1(J_i)$} $\simeq$ \unsure{Cl.|cl.} de Hodge}}}

Moyennant un certain nombre de choses non prouvées (qui risquent d'ailleurs d'être vaches) le théorème de positivité de Hodge se ramène à l'énoncé suivant, que je ne sais pas prouver, et que j'ai envie de te communiquer car il est bien terre à terre. Moralement, il s'agit de prouver que dans le cas où $\dim X=2m-1$, l'autodualité de $J^m$
% slanted
\margin{\handwritten{\unsure{canulé|annulé|Cancelé} tel quel, il faut tenir compte de la décomposition primitive\ldots\ldots}}
(qui s'exprime par une classe de correspondance divisorielle sur $J^m\mathbin{\mathrm{x}}J^m$ \deleted{anti}\uline{symétrique}, donc provenant - du moins modulo le facteur 2-
